\section*{Capítulo II, Problema IX}

\textbf{Problema:}

9. Encuentra un vector perpendicular a $(1,2,-3)$ y $(2,-1,3)$, y otro vector perpendicular a $(-1,3,2)$ y $(2,1,1)$.

\textbf{Solución:}

Un vector perpendicular a dos vectores en $\mathbb{R}^3$ se obtiene con su producto cruz.

\begin{itemize}
    \item Para $(1,2,-3)$ y $(2,-1,3)$:
    \[
        (1,2,-3) \times (2,-1,3)
        = \bigl(2\cdot 3 - (-3)(-1),\; (-3)\cdot 2 - 1\cdot 3,\; 1\cdot (-1) - 2\cdot 2\bigr).
    \]
    Por tanto,
    \[
        (1,2,-3) \times (2,-1,3) = (3,-9,-5).
    \]
    Verificamos:
    \[
        (3,-9,-5) \cdot (1,2,-3) = 3 - 18 + 15 = 0,
    \]
    \[
        (3,-9,-5) \cdot (2,-1,3) = 6 + 9 - 15 = 0.
    \]

    \item Para $(-1,3,2)$ y $(2,1,1)$:
    \[
        (-1,3,2) \times (2,1,1)
        = \bigl(3\cdot 1 - 2\cdot 1,\; 2\cdot 2 - (-1)\cdot 1,\; (-1)\cdot 1 - 3\cdot 2\bigr).
    \]
    Así obtenemos
    \[
        (-1,3,2) \times (2,1,1) = (1,5,-7).
    \]
    En efecto,
    \[
        (1,5,-7) \cdot (-1,3,2) = -1 + 15 - 14 = 0,
    \]
    \[
        (1,5,-7) \cdot (2,1,1) = 2 + 5 - 7 = 0.
    \]
\end{itemize}

Por consiguiente, una respuesta posible es
\[
    (3,-9,-5) \quad \text{y} \quad (1,5,-7).
\]